% Chapter 3: Supervised Learning
\chapter{Supervised Learning}

\section{Problem Formulation}
The prediction task is formulated as binary classification:
\begin{itemize}
    \item \textbf{Class 1:} Fire detected at location
    \item \textbf{Class 0:} No fire detected at location
\end{itemize}

\textbf{Class balancing decision:} The original dataset was imbalanced. We applied undersampling to balance classes, as this approach is simpler than oversampling techniques like SMOTE and sufficient given our data volume.

\section{Feature Engineering}

\subsection*{Feature Selection}
We applied multiple selection techniques:
\begin{enumerate}
    \item \textbf{Correlation analysis:} Features with correlation $> 0.95$ removed to reduce multicollinearity
    \item \textbf{Variance threshold:} Low-variance features removed
    \item \textbf{Recursive Feature Elimination:} Using Random Forest importance
\end{enumerate}

\textbf{Decision rationale:} High correlation between features causes instability in model coefficients and makes interpretation difficult. We prioritized keeping features with clear physical interpretation.

\subsection*{Derived Features}
Interaction features created based on domain knowledge:
\begin{itemize}
    \item \texttt{lat\_long\_interaction}: Captures regional patterns
    \item \texttt{elevation\_slope\_interaction}: Terrain effects on fire spread
    \item \texttt{drought\_index}: Combined climate stress indicator
\end{itemize}

\begin{figure}[H]
    \centering
    \fbox{\parbox{0.85\textwidth}{\centering\vspace{3cm}\textbf{[PLACEHOLDER: Feature Correlation Heatmap]}\vspace{3cm}}}
    \caption{Correlation matrix showing relationships between features}
    \label{fig:feature_correlation}
\end{figure}

\subsection*{Feature Scaling}
Scaling applied based on algorithm requirements:
\begin{itemize}
    \item \textbf{KNN:} StandardScaler (z-score)---essential for distance-based algorithms
    \item \textbf{Decision Tree/Random Forest:} No scaling---tree-based splits are invariant to monotonic transformations
\end{itemize}

\section{K-Nearest Neighbors (KNN)}

\subsection*{Algorithm Choice Rationale}
KNN selected as baseline model because: (1) simple and interpretable, (2) makes no assumptions about data distribution, (3) good for detecting local patterns.

\subsection*{Hyperparameter Tuning}
Optimal $k$ selected via 5-fold cross-validation:
\begin{itemize}
    \item Search range: $k \in \{3, 5, 7, 9, 11, 15, 21, 31\}$
    \item Metric: F1-score (handles class imbalance)
    \item Distance: Euclidean on standardized features
\end{itemize}

\begin{figure}[H]
    \centering
    \fbox{\parbox{0.85\textwidth}{\centering\vspace{3cm}\textbf{[PLACEHOLDER: KNN Cross-Validation Results]}\vspace{3cm}}}
    \caption{Cross-validation performance for different $k$ values}
    \label{fig:knn_cv}
\end{figure}

\subsection*{Results and Analysis}
\begin{itemize}
    \item \textbf{Strengths:} Good local pattern detection, interpretable
    \item \textbf{Limitations:} Computationally expensive at prediction time, sensitive to irrelevant features
\end{itemize}

\section{Decision Tree}

\subsection*{Algorithm Choice Rationale}
Decision Tree provides: (1) interpretable decision rules, (2) handles non-linear relationships, (3) feature importance ranking.

\subsection*{Model Configuration}
\begin{itemize}
    \item Criterion: Gini impurity
    \item Hyperparameters tuned: \texttt{max\_depth}, \texttt{min\_samples\_split}, \texttt{min\_samples\_leaf}
    \item Validation: 5-fold stratified cross-validation
\end{itemize}

Gini impurity for node splitting:
\begin{equation}
    G = 1 - \sum_{i=1}^{C} p_i^2
\end{equation}

\begin{figure}[H]
    \centering
    \fbox{\parbox{0.85\textwidth}{\centering\vspace{3cm}\textbf{[PLACEHOLDER: Decision Tree Visualization]}\vspace{3cm}}}
    \caption{Trained Decision Tree structure (limited depth)}
    \label{fig:dt_viz}
\end{figure}

\subsection*{Key Decision Rules}
The tree reveals interpretable patterns:
\begin{itemize}
    \item Primary splits on climate features (drought index, summer temperature)
    \item Coastal proximity as important early split
    \item Elevation for fine-grained distinctions
\end{itemize}

\section{Random Forest}

\subsection*{Algorithm Choice Rationale}
Random Forest selected as primary model because: (1) ensemble reduces overfitting, (2) robust feature importance, (3) handles high-dimensional data well.

\subsection*{Model Configuration}
\begin{itemize}
    \item Bootstrap aggregation with random feature selection
    \item Hyperparameters: \texttt{n\_estimators}, \texttt{max\_depth}, \texttt{max\_features}
    \item Out-of-bag (OOB) score for additional validation
\end{itemize}

\subsection*{Results}
\begin{table}[H]
\centering
\caption{Random Forest Performance on Test Set}
\label{tab:rf_performance}
\begin{tabular}{|l|c|}
\hline
\textbf{Metric} & \textbf{Value} \\
\hline
Accuracy & 0.959 \\
Precision & 0.973 \\
Recall & 0.889 \\
F1-Score & 0.929 \\
AUC-ROC & 0.989 \\
\hline
\end{tabular}
\end{table}

\begin{figure}[H]
    \centering
    \fbox{\parbox{0.85\textwidth}{\centering\vspace{3cm}\textbf{[PLACEHOLDER: Random Forest Confusion Matrix and ROC Curve]}\vspace{3cm}}}
    \caption{Confusion matrix (1766 TN, 680 TP, 19 FP, 85 FN) and ROC curve}
    \label{fig:rf_results}
\end{figure}

\subsection*{Feature Importance}
\begin{table}[H]
\centering
\caption{Top 10 Feature Importance Scores}
\label{tab:rf_importance}
\begin{tabular}{|l|c|l|}
\hline
\textbf{Feature} & \textbf{Importance} & \textbf{Interpretation} \\
\hline
coastal\_proximity & 0.114 & Coastal Mediterranean fire concentration \\
lat\_long\_interaction & 0.067 & Regional geographic patterns \\
elevation & 0.060 & Topographic influence on fires \\
drought\_index & 0.058 & Climate stress indicator \\
prec\_winter & 0.054 & Antecedent moisture conditions \\
elevation\_slope\_interaction & 0.044 & Terrain fire spread effects \\
tmax\_summer & 0.043 & Peak heat conditions \\
prec\_autumn & 0.040 & Pre-fire season moisture \\
total\_annual\_precip & 0.039 & Overall moisture availability \\
terrain\_complexity & 0.039 & Rugged terrain effects \\
\hline
\end{tabular}
\end{table}

\begin{figure}[H]
    \centering
    \fbox{\parbox{0.85\textwidth}{\centering\vspace{3cm}\textbf{[PLACEHOLDER: Feature Importance Bar Chart]}\vspace{3cm}}}
    \caption{Complete feature importance ranking}
    \label{fig:rf_importance}
\end{figure}

\section{Model Comparison}

\begin{table}[H]
\centering
\caption{Supervised Learning Model Comparison}
\label{tab:model_comparison}
\begin{tabular}{|l|c|c|c|c|c|}
\hline
\textbf{Model} & \textbf{Accuracy} & \textbf{Precision} & \textbf{Recall} & \textbf{F1} & \textbf{AUC} \\
\hline
KNN & -- & -- & -- & -- & -- \\
Decision Tree & -- & -- & -- & -- & -- \\
Random Forest & 0.959 & 0.973 & 0.889 & 0.929 & 0.989 \\
\hline
\end{tabular}
\end{table}

\textbf{Model selection decision:} Random Forest achieves the best performance across all metrics. Its ensemble nature reduces overfitting while providing interpretable feature importance. The high precision (0.973) minimizes false alarms, while good recall (0.889) captures most fire events.

\section{Summary}
Key findings from supervised learning:
\begin{enumerate}
    \item \textbf{Best model:} Random Forest with 95.9\% accuracy and 0.989 AUC-ROC
    \item \textbf{Most important predictor:} Coastal proximity---fires concentrate in coastal Mediterranean zones
    \item \textbf{Climate importance:} Drought index and precipitation patterns are critical predictors
    \item \textbf{Practical application:} Model can identify high-risk areas for targeted fire prevention
\end{enumerate}

Detailed hyperparameter search results are provided in Appendix~\ref{app:hyperparams}.
